\section{A Spectral Sequence for Extensions}\label{sec:ess}
In this paper, we assume that all symbols $X,Y,Z,W,\dots$ for spectra represent 2-completed connective spectra of finite type. In particular, their $\HF$-Adams spectral sequences converge strongly.

Consider a map $f: X\to Y$ between two spectra and we can construct the following chain complex
$$0\to \pi_*X\fto{\pi_*f} \pi_*Y\to 0$$
whose homology is 
$$\Ker(\pi_*f)\oplus \Coker(\pi_*f).$$
We encounter $f$-extension problems when analyzing the homotopy groups through the $E_\infty$-pages of the $\HF$-based Adams spectral sequences. To address these extensions, we introduce the following spectral sequence.



\begin{definition}\label{def:ess}
  For a map $f: X\to Y$, we define an $f$-extension spectral sequence (denoted $f$-ESS) as the spectral sequence derived from filtering the above chain complex according to the Adams filtration on $\pi_*$. The $E_0$-page of this spectral sequence is isomorphic to the direct sum of $E_\infty$-pages of the Adams spectral sequences of $X$ and $Y$: 
  $$\ESS{f}_0^{s,t} \cong E_\infty^{s,t}(X) \oplus E_\infty^{s,t}(Y) \Longrightarrow \Ker(\pi_*f)\oplus \Coker(\pi_*f)$$
  where the $d_0$-differential is induced by $f$ on the $E_\infty$-pages.
  
  Differentials in this spectral sequence has the form
  $$d_n^{f}: \ESS{f}_n^{s,t}\to \ESS{f}_n^{s+n, t+n}.$$
\end{definition}

To prevent confusion between differentials $d_n$ in the Adams spectral sequence of a spectrum $X$ and those in the $f$-extension spectral sequence, we denote the latter with a superscript, writing them as $d_n^{f}$. 

In the remainder of the paper, we denote by $\AF(f)$ the Adams filtration of a map $f$. Recall that $\AF(f) \ge k$ if and only if there exist $k$ composable maps $f_1, \ldots, f_k$, each inducing a trivial map in $\HF$-homology, such that $f$ is homotopic to the composition $f_k \circ \cdots \circ f_1$ (see \cite[Theorem~2.2.14]{RavenelGreenbook} for example).

\begin{notation}\label{nota:ess-notation}
    \uses{def:ess}
    Let $\ZESS{f}^{s,t}_n(X)\subset E^{s,t}_\infty(X)$ be the subgroup consisting of elements for which the $f$-extension differentials $d^f_0,\dots,d^f_n$ vanish.
    Let $\BESS{f}^{s,t}_n(Y)\subset E^{s,t}_\infty(Y)$ be the subgroup generated by the sum of images of the $f$-extension differentials $d^f_0,\dots,d^f_n$. We define $\BESS{f}^{s,t}_{-1}(Y)=0$ and $\ZESS{f}^{s,t}_{-1}(X)=E^{s,t}_\infty(X)$.\\
    
    By definition we have
    $$\ESS{f}^{s,t}_n\iso \ZESS{f}^{s,t}_{n-1}(X)\oplus \big(E^{s,t}_\infty(Y)/\BESS{f}^{s,t}_{n-1}(Y)\big).$$
    
    A nonzero $d^f_n$ differential takes the form
    $$d^f_n: \ZESS{f}^{s,t}_{n-1}(X)\to E^{s+n,t+n}_\infty(Y)/\BESS{f}^{s+n,t+n}_{n-1}(Y).$$
     When we write 
    $$d^f_n(x)=y$$
    for $x\in E_\infty^{s,t}(X)$ and $y\in E_\infty^{s+n,t+n}(Y)$, it is understood that $x$ belongs to $\ZESS{f}^{s,t}_{n-1}(X)$, and 
    $$d^f_n(x)=y+\BESS{f}^{s+n,t+n}_{n-1}(Y).$$
    Furthermore, we may sometimes write
    $$d^f_n(x)\equiv y\mod M$$
    for some subgroup $M\subset E_\infty^{s+n,t+n}(Y)$. This means that 
    $$d^f_n(x)=y+m+\BESS{f}^{s+n,t+n}_{n-1}(Y)$$
    for some $m\in M$.
\end{notation}

\begin{definition}\label{def:f-extension}
    \uses{def:ess,nota:ess-notation}
    We say there is an $f$-extension from $x\in E_\infty^{s,t}(X)$ to $y\in E_\infty^{s+n,t+n}(Y)$ if $d_n^{f}(x)=y$ in the $f$-ESS. We say this $f$-extension is $\textbf{essential}$ if $y$ is nontrivial in the $E_n$-page of the $f$-ESS, or equivalently $$y\notin \BESS{f}^{s+n,t+n}_{n-1}(Y).$$
    Otherwise we refer to it as inessential.
\end{definition}

\begin{notation}
    For $x\in E_\infty^{s,t}(X)$, we denote by $\{x\}$  the set of all classes in $\pi_{t-s}X$ that are detected by $x$. We use $[x]$ to refer to a specific class or a general class in $\{x\}$, with the choice understood from the context.
\end{notation}


\begin{proposition}\label{prop:f-ext-characterization}
    \uses{def:f-extension}
Consider $f: X\to Y$, $x\in E_\infty^{s,t}(X)$, $y\in E_\infty^{s+n,t+n}(Y)$ and $y'\in E_\infty^{s+m,t+m}(Y)$ for $m,n\ge 0$.
\begin{enumerate}
    \item There is an $f$-extension from $x$ to $y$, i.e., $d_n^{f}(x)=y$ in the $f$-ESS if and only if there is a class $[x]\in \{x\}$ such that $f[x]\in \{y\}$.
    \item An $f$-extension from $x$ to $y$ is inessential, i.e., $y$ is trivial in the $E_n$-page of the $f$-ESS if and only if there exists an element $x'\in E_\infty^{s+a,t+a}(X)$ for $0<a\le n$ such that we have an essential differential $d_{n-a}^{f}(x')=y$. Equivalently, there exists a class $[x'] \in \{x'\} \subset \pi_*X$ with AF$(x')>$AF$(x)$, such that $f[x'] \in \{y\}$.
    
    \item Suppose we have an $f$-extension from $x$ to both $y$ and $y'$.
    \begin{enumerate}
        \item If $m=n$, then $y-y'\in \BESS{f}^{s+n,t+n}_{n-1}(Y)$ and there exists an element $x'\in E_\infty^{s+a,t+a}(X)$ for $0<a\le n$ such that we have an essential differential 
        $$d_{n-a}^{f}(x') = y-y'$$
        in $f$-ESS. Equivalently, there exists a class $[x'] \in \{x'\} \subset \pi_*X$ with AF$(x')>$AF$(x)$, such that $f[x'] \in \{y-y'\}$. 
        \item If $m>n$, then the $f$-extension from $x$ to $y$ is inessential.
     \end{enumerate}
\end{enumerate}
\end{proposition}

\begin{proof}
    Part (1) is straightforward from the setup of the extension spectral sequence and Definition~\ref{def:f-extension}. Parts (2) and (3) follow from Part (1).
\end{proof}

\begin{corollary}\label{cor:af-vanishing}
    \uses{def:ess}
    If the Adams filtration of $f$ is $k$, then
    $d^f_i=0$ for $i<k$.
\end{corollary}

\begin{proof}
    When $k>0$, it is clear that $d^f_0=0$, which means that for any $s$ and $[x]\in F_s\pi_*(X)$, we have $f([x])\in F_{s+1}\pi_*(Y)$. Since $\AF(f)=k$, we know that there exist $k$ maps
    $$X=X_0\fto{f_1}X_1\fto{f_2}\cdots\fto{f_k}X_k=Y$$
    such that each $f_i$ induces trivial map in $\HF$-homology (hence $\AF(f_i)>0$) and $f$ is homotopic to the composition $f_k\circ \cdots\circ f_1.$
    Then for any $s$ and $[x]\in F_s\pi_*(X)$, inductively we have
    $$f_1([x])\in F_{s+1}\pi_*(X_1)$$
    $$f_2\circ f_1([x])\in F_{s+2}\pi_*(X_2)$$
    $$\cdots$$
    $$f([x])=f_k\circ\cdots\circ f_1([x])\in F_{s+k}\pi_*(X_k)=F_{s+k}\pi_*(Y).$$
    By Proposition \ref{prop:f-ext-characterization}.(1) this means that $d_i^fx=0$ for any $x\in E_\infty^{*,*}(X)$ and $i<k$.
\end{proof}

\begin{definition} \label{def:crossing}
\uses{def:ess}
\begin{enumerate}
    \item For $x\in E_\infty^{s,t}(X)$, we say that the an $f$-extension $d_n^{f}(x)=y$ has a crossing that hits Adams filtration $p$ for some $p \leq s+n$, if there exists an element $x^\prime\in E_\infty^{s+a,t+a}(X)$ with $a>0$ and an $f$-extension $d_m^{f}(x^\prime)=y^\prime$ for $0\neq y^\prime\in E^{s+a+m, t+a+m}_\infty(Y)$ such that
    $$p = \textup{AF}(y^\prime) = s+a+m\le \textup{AF}(y) = s+n.$$
    \item We say that an $f$-extension $d_n^{f}(x)=y$ has no crossing that hits the range AF $\ge p$ if there does not exist an element $x^\prime\in E_\infty^{s+a,t+a}(X)$ with $a>0$ and an $f$-extension $d_m^{f}(x^\prime)=y^\prime$ for $0\neq y^\prime\in E^{s+a+m, t+a+m}_\infty(Y)$ such that
    $$p\le \textup{AF}(y^\prime) = s+a+m\le \textup{AF}(y) = s+n.$$
    \item We say that an $f$-extension $d_n^{f}(x)=y$ has no crossing if it has no crossing that hits the range AF$\ge \AF(x) + 1$.
\end{enumerate}
\end{definition}


\begin{figure}[h]
\centering
\scalebox{0.9}{\begin{tikzpicture}
    % Draw grid
    \fill[llgray] (2-0.5, 5+0.5) rectangle ++(1, -3);
    \tikzgrid{0}{1}{0}{6}
    \tikzgrid{2}{3}{0}{6}
    \draw[thick] (2-0.5, 5+0.5) rectangle ++(1, -3);

    % Draw bullets
    \coordinate (x) at (0, 0);
    \coordinate (y) at (2, 5);
    \coordinate (x1) at (0, 1);
    \coordinate (y1) at (2, 3);
    
    \fill (x) circle (0.1);
    \fill (y) circle (0.1);
    \fill (x1) circle (0.1);
    \fill (y1) circle (0.1);

    % Draw differentials
    \drawarrow (x) node[left] {$x~$} -- (y) node[right] {$y$};
    \drawarrow (x1) node[left] {$x^\prime$} -- (y1) node[right] {$y^\prime$};
    
    \node[anchor=south east] at (1.5, 3.5) {$d^f_n$};
    \node[anchor=north west] at (1.2, 2.5) {$d^f_m$};

    \node[anchor=east] at (-1, 0) {$s$};
    \node[anchor=east] at (-1, 5) {$s+n$};
    \node[anchor=east] at (-1, 1) {$s+a$};
    \node[anchor=east] at (-1, 3) {$p=s+a+m$};

    
    \node at (0, -1.0) {$E_\infty(X)$};
    \node at (2, -1.0) {$E_\infty(Y)$};
\end{tikzpicture}}
\caption{A crossing of $d_n^f$ that hits Adams filtration $p$}
\end{figure}
   
\begin{remark}   

The following statements are equivalent.
\begin{itemize}
    \item An $f$-extension $d_n^{f}(x)=y$ has no crossing that hits the range AF$\ge p$. 
    \item If for any $a>0$ there is an $f$-extension from $x^\prime\in E_\infty^{s+a,t+a}(X)$ to some nontrivial $y^\prime$, then AF$(y^\prime)<p$ or AF$(y^\prime)>$AF$(y)$.
\end{itemize} 
\end{remark}
    
\begin{remark}   
We do not require the crossings in Definition~\ref{def:crossing}~(1)(2) to be essential. A crossing implies the existence of an essential crossing, possibly a shorter one.
\end{remark}


\begin{proposition} \label{prop:no-crossing-iff}
    \uses{prop:f-ext-characterization,def:crossing}
    An $f$-extension from $x$ to $y$ has no crossing in the range $\AF\ge p$ if and only if for all $[x]\in \{x\}$ such that $\AF(f[x])\ge p$ we have $f[x]\in \{y\}$. In particular, we have
    \begin{enumerate}
        \item An $f$-extension from $x$ to $y$ has no crossing if and only if for all $[x]\in \{x\}$ we have $f[x]\in \{y\}$.
        \item An $f$-extension $d_n^{f}(x)=0$ has no crossing if and only if for all $[x]\in \{x\}$,
        $$\AF(f[x]) >\AF(x)+n.$$
        \item If $\AF(f)=n$, then all $d^f_n$-differentials have no crossing in the $f$-ESS.
    \end{enumerate}
\end{proposition}

\begin{proof}
    First we prove the only if part. Assume by contradiction that there exists $[x]\in \{x\}$ such that $f[x]=[y']\notin \{y\}$ is detected by $y'\in E_\infty(Y)$ with $\AF(y')\ge p$. Then there is an $f$ extension from $x$ to $y'$. By Proposition \ref{prop:f-ext-characterization}~(3), $y'$ (if $\AF(y)>\AF(y')$)  or $y-y'$ (if $\AF(y)=\AF(y')$) or $y$ (if $\AF(y)<\AF(y')$) should be hit by a shorter $d^f$ differential. This shorter differential is a crossing of the $f$-extension from $x$ to $y$ that hits $\AF\ge p$, which is a contradiction. Now we prove the if part. Assume that the $f$-extension has a crossing from $x'$ to $y'$ with $p\le \AF(y')\le \AF(y)$. Then there exists $[x']\in \{x'\}$ such that $f([x'])\in \{y'\}$. Note that $[x]+[x']\in \{x\}$ since $\AF(x')>\AF(x)$, we have
    $$f([x]+[x'])=f([x])+f([x'])=[y]+[y']\in\begin{cases}
        \{y'\} & \text{ if } \AF(y')<\AF(y)\\
        \{y+y'\} & \text{ if } \AF(y')=\AF(y).\\
    \end{cases}$$
    However, by assumption we should have $f([x]+[x'])\in \{y\}$ and it is contradictory to the equation above in either case.
    
    Hence the main statement is true.  Part (1) (2) are special cases of the main statement and Part (3) holds for degree reasons.
\end{proof}





\begin{theorem}\label{thm:four-spectra}
    \uses{prop:no-crossing-iff}
    Consider a homotopy commutative diagram of spectra
    $$\xymatrix{
    X \ar[r]^{f}\ar[d]_{p} & Y\ar[d]^{q}\\
    Z \ar[r]_{g} & W
    }$$
    Suppose $m,n,l\ge 0$, $0<k\le m+l-n$,
    $$\begin{array}{ll}
        x\in E^{s,t}_\infty(X) & y\in E^{s+n,t+n}_\infty(Y),\\
        z\in E^{s+m,t+m}_\infty(Z) & w\in E^{s+m+l,t+m+l}_\infty(W),
    \end{array}$$
    and
    \begin{enumerate}
        \item $d^{f}_{n}(x)=y$,
        \item $d^{p}_{m}(x)=z$,
        \item the differential in (1) or the differential in (2) has no crossing,
        \item $d^{g}_{l}(z)=w$, and has no crossing that hits the range AF $\ge s+n+k$,
        \item $d^{q}_{k-1} y=0$, and has no crossing.
    \end{enumerate}
    Then $d^{q}_{m+l-n}(y)=w$. (See Figure \ref{fig:four-spectra}.)
\end{theorem}

\begin{proof}
    First we find a representative $[x]$ of $x$ such that
    \begin{equation}\label{eq:p-representative}
        p[x]\in \{z\}
    \end{equation}
    and
    \begin{equation}\label{eq:f-representative}
        f[x]\in \{y\}.
    \end{equation}
    If (1) has no crossing, by (2) we can pick $[x]\in \{x\}$ such that (\ref{eq:p-representative}) holds;
    If (2) has no crossing, by (1) we can pick $[x]\in \{x\}$ such that  (\ref{eq:f-representative}) holds.
    In both cases the other equation also holds because of the no crossing condition.
    
    By (5), (\ref{eq:f-representative}) and Proposition~\ref{prop:no-crossing-iff}~(2) we know that $qf[x]$ has AF $\ge s+n+k$. Since the diagram on the Adams $E_\infty$-pages commutes, we have
    $$AF(gp[x])=AF(qf[x])\ge s+n+k.$$
    Combining with (4), (\ref{eq:p-representative}) and Proposition~\ref{prop:no-crossing-iff} we have
    $$gp[x]\in \{w\}.$$
    Therefore, $qf[x]\in \{w\}$ and by (\ref{eq:f-representative}) there is a $q$-extension from $y$ to $w$.
\end{proof}

\begin{figure}
\centering
\scalebox{0.9}{\begin{tikzpicture}[node/.style={circle,fill=black!0.1}]
    \def\ymax{6}
    \def\ymaxp{7}
    % Define colors
    \definecolor{lightgray}{RGB}{200,200,200}
    
    \fill[llgray] (6-0.5, 6+0.5) rectangle ++(1, -2);
    % Draw grid
    \foreach \x in {0,...,7} {
        \draw[lightgray] (\x-0.5, -0.5) -- (\x-0.5, \ymax+0.5);
    }
    \foreach \x in {0,2,4,6} {
        \foreach \y in {0,...,\ymaxp} {
            \draw[lightgray] (\x-0.5, \y-0.5) -- (\x+0.5, \y-0.5);
        }
    }
    \draw[thick] (6-0.5, 6+0.5) rectangle ++(1, -2);

    % Draw bullets
    \coordinate (x) at (0, 0);
    \coordinate (y) at (4, 1);
    \coordinate (z) at (2, 4);
    \coordinate (w) at (6, 6);
    
    \fill (x) circle (0.1) node[below left] {$x$};
    \fill (y) circle (0.1) node[below right] {$y$};
    \fill (z) circle (0.1) node[above left] {$z$};
    \fill (w) circle (0.1) node[above right] {$w$};

    % Draw differentials
    \drawarrow (x) -- (y) node[midway,below] {$d^f_n$};
    \drawarrow (x) -- (z) node[midway,left] {$d^p_m$};
    \draw[dashed, -{Stealth[length=1.5mm]}, shorten >=(0.08cm)] (y) -- (6,6-1) node[midway,right] {$d^q_{\ge k}$};
    \drawarrow (z) -- (w) node[midway,above] {$d^g_{l}$};

    \node[anchor=east] (axis1) at (-.7, 0) {$s$};
    \node[anchor=east] (axis2) at (-.7, 1) {$s+n$};
    \node[anchor=east] (axis3) at (-.7, 4) {$s+m$};
    \node[anchor=east] (axis4) at (-.7, 5) {$s+n+k$};
    \node[anchor=east] (axis5) at (-.7, 6) {$s+m+l$};

    \node (EX) at (0, -1.0) {$E_\infty(X)$};
    \node (EY) at (4, -1.0) {$E_\infty(Y)$};
    \node (EZ) at (2, -1.0) {$E_\infty(Z)$};
    \node (EW) at (6, -1.0) {$E_\infty(W)$};
\end{tikzpicture}}
\caption{A demonstration of Theorem \ref{thm:four-spectra}}\label{fig:four-spectra}
\end{figure}


\begin{corollary} \label{cor:four-spectra}
    \uses{thm:four-spectra}
    Consider a homotopy commutative diagram of spectra
    $$\xymatrix{
    X \ar[r]^{f}\ar[d]_{p} & Y\ar[d]^{q}\\
    Z \ar[r]_{g} & W
    }$$
    Suppose $m,n,l\ge 0$,
    $$\begin{array}{ll}
        x\in E^{s,t}_\infty(X) & y\in E^{s+n,t+n}_\infty(Y),\\
        z\in E^{s+m,t+m}_\infty(Z) & w\in E^{s+m+l,t+m+l}_\infty(W),
    \end{array}$$
    and
    \begin{enumerate}
        \item $d^{f}_{n}(x)=y$,
        \item $d^{p}_{m}(x)=z$,
        \item the differential in (1) or the differential in (2) has no crossing,
        \item $d^{g}_{l}(z)=w$, and has no crossing.
    \end{enumerate}
    Then $d^{q}_{m+l-n}(y)=w$.
\end{corollary}

\begin{corollary} \label{cor:triangle-map}
    \uses{cor:four-spectra}
    Consider a homotopy commutative diagram of spectra
    $$\xymatrix{
    X \ar[r]^{f}\ar[dr]_{p} & Y\ar[d]^{q}\\
    & Z
    }$$
    Suppose $n,m\ge 0$, 
    $$x\in E^{s,t}_\infty(X), ~y\in E^{s+n,t+n}_\infty(Y), ~z\in E^{s+m,t+m}_\infty(Z)$$
    and
    \begin{enumerate}
        \item $d^{f}_{n}(x)=y$,
        \item $d^{p}_{m}(x)=z$,
        \item the differential in (1) or the differential in (2) has no crossing.
    \end{enumerate}
    Then $d^{q}_{m-n}(y)=z$.
\end{corollary}

\begin{proof}
This is a special case of Corollary \ref{cor:four-spectra} when $Z=W$ and $g=id_Z$.
\end{proof}

\begin{corollary} \label{cor:composition-ext}
    \uses{cor:four-spectra}
    Consider a homotopy commutative diagram of spectra
    $$\xymatrix{
    X \ar[d]_{p}\ar[dr]^{q} & \\
    Z\ar[r]_{g} & W
    }$$
    Suppose $n,m\ge 0$, 
    $$x\in E^{s,t}_\infty(X), ~z\in E^{s+m,t+m}_\infty(Y), ~w\in E^{s+m+l,t+m+l}_\infty(W)$$
    and
    \begin{enumerate}
        \item $d^{p}_{m}(x)=z$,
        \item $d^{g}_{l}(z)=w$, and has no crossing.
    \end{enumerate}
    Then $d^{q}_{m+l}(x)=w$.
\end{corollary}

\begin{proof}
This is a special case of Corollary \ref{cor:four-spectra} when $X=Y$ and $f=id_X$.
\end{proof}

\begin{corollary}\label{cor:ess-naturality}
    \uses{def:ess}
    Consider a homotopy commutative diagram of spectra
    $$\xymatrix{
    X \ar[r]^{f}\ar[d]_{p} & Y\ar[d]^{q}\\
    Z \ar[r]_{g} & W
    }.$$
    Assume $\ESS{p}_0^{*,*}=\ESS{p}_r^{*,*}$ and $\ESS{q}_0^{*,*}=\ESS{q}_r^{*,*}$ for some $r\ge 0$. Then 
    $$(d^p_r, d^q_r): E_\infty^{*,*}(X)\oplus E_\infty^{*,*}(Y)\to E_\infty^{*,*}(Z)\oplus E_\infty^{*,*}(W)$$
    induces a map from the $f$-ESS to the $g$-ESS.
\end{corollary}

\begin{proof}
    Since $\ESS{p}_0^{*,*}=\ESS{p}_r^{*,*}$ and $\ESS{q}_0^{*,*}=\ESS{q}_r^{*,*}$, we know that $d^p_r$, $d^q_r$ have no crossings. By Corollary \ref{cor:four-spectra}, we have $d^g_0\circ d^p_r = d^q_r\circ d^f_0$ and therefore $(d^p_r, d^q_r)$ induces a map from $\ESS{f}_1^{*,*}$ to $\ESS{g}_1^{*,*}$. Inductively we can apply Corollary \ref{cor:four-spectra} again and show that $(d^p_r, d^q_r)$ induces a map from $\ESS{f}_n^{*,*}$ to $\ESS{g}_n^{*,*}$ and $d^g_n\circ d^p_r = d^q_r\circ d^f_n$ for all $n$.
\end{proof}

\begin{corollary}\label{cor:composition-zero}
    \uses{cor:four-spectra,prop:f-ext-characterization}
    If the composition $g\circ f$ of two maps $f:X\to Y$, $g:Y\to Z$ is trivial, and $d_n^f(x)=y$ for $x\in E_\infty^{s,t}(X)$ and $y\in E_\infty^{s+n,t+n}(Y)$, then $y$ is a permanent cycle in the $g$-ESS, i.e., we have
    $d^g_m(y)=0$ for all $m\ge 0$.
\end{corollary}

\begin{proof}
    This follows immediately from the following commutative diagram
    $$\xymatrix{
    X \ar[r]^{f}\ar[d] & Y\ar[d]^{g}\\
    0 \ar[r] & Z
    }$$
    and Corollary \ref{cor:four-spectra}.

    There is a second short proof using homotopy groups. By Proposition \ref{prop:f-ext-characterization}, we know that there exists $[x]\in \{x\}$ such that $f([x])\in \{y\}$. Let $[y]=f([x])\in \{y\}$ and we have $g([y])=g(f([x]))=0$. Hence $y$ is a permanent $d^g$ cycle.
\end{proof}

\begin{proposition}\label{prop:ess-exactness}
    \uses{prop:f-ext-characterization,cor:composition-zero}
    Suppose that the sequence of spectra
    $$X\fto{f}Y\fto{g}Z$$
    induces a sequence of homotopy groups
    $$\pi_*X\fto{\pi_*f}\pi_*Y\fto{\pi_*g}\pi_*Z$$
    that is exact in $\pi_*Y$. Then all permanent cycles in $E^{*,*}_\infty(Y)$ of the $g$-ESS are boundaries in the $f$-ESS.
\end{proposition}

\begin{proof}
    If we consider the following sequence
    \begin{equation}\label{eq:exact-chain}
        0\to \pi_*X\fto{\pi_*f}\pi_*Y\fto{\pi_*g}\pi_*Z\to 0
    \end{equation}
    and treat it as a chain complex filtered by the Adams filtration, we obtain an extension spectral sequence comprising $d^f$ and $d^g$ differentials. The abutment of this spectral sequence, when projected onto the $Y$ component, is zero. Therefor, all permanent $d^g$-cycles are killed by $d^f$-differentials.

    Again we offer another proof via homotopy groups. Assume that $y\in E^{*,*}_\infty(Y)$ is a permanent $d^g$ cycle. By the convergence of ESS we know that it detects an $g$-torsion $[y]\in \{y\}$ with $g([y])=0$. Hence there exists $[x]\in \pi_*X$ such that $f([x])=[y]$ by exactness. Assume that $[x]$ is detected by $x\in E^{*,*}_\infty(X)$. Then there exists an $f$ extension from $x$ to $y$.
\end{proof}

Corollary \ref{cor:composition-zero} and Proposition \ref{prop:ess-exactness} are particularly useful because they provide numerous potential applications whenever we have a cofiber sequence of the form
$$X\fto{f}Y\fto{g}Cf\fto{h}\Sus X$$
These results can also be implemented in a computer program to facilitate their application.

